% Build with XeLaTeX; this file is self-contained.
\documentclass[aspectratio=169,10pt]{beamer}
\usepackage{fontspec}
\usepackage[russian]{babel}
\setsansfont{Arial}
\setmainfont{Arial}
\usefonttheme{professionalfonts}
\usepackage{amsmath,pgfplots,pgfplotstable}
\pgfplotsset{compat=1.18}

\definecolor{Ink}{HTML}{222222}
\definecolor{Accent}{HTML}{176B75}
\definecolor{Muted}{HTML}{A9B2B8}
\definecolor{Subtle}{HTML}{555555}
\setbeamercolor{normal text}{fg=Ink,bg=white}
\setbeamercolor{frametitle}{fg=Ink,bg=white}
\setbeamercolor{structure}{fg=Accent}
\setbeamerfont{frametitle}{size=\fontsize{17}{20},series=\bfseries}
\setbeamertemplate{navigation symbols}{}
\setbeamertemplate{footline}{}
\setbeamertemplate{frametitle}{\vspace{2.5mm}\insertframetitle\par\vspace{.8mm}}
\setbeamersize{text margin left=7mm,text margin right=7mm}
\setlength{\parskip}{1mm}
\hypersetup{
  pdftitle={Перенос VQDM на DiT: важны средние шаги},
  pdfauthor={Егор Подъяблонский},
  pdfsubject={Однослайдовый доклад о переносе векторной квантизации на DiT},
  pdfcreator={XeLaTeX},
  hidelinks
}

% Source: ../article_assets/results.json, temporal_windows_256.
% All bounds are paired-bootstrap 95% intervals; the first crosses zero.
\pgfplotstableread{
position gain low high
0 2.368293 -0.648391 4.993902
1 20.627143 16.706798 24.487416
2 36.488425 31.108123 41.911975
3 12.799602 9.918571 15.986773
4 1.089730 0.509793 1.694293
}\windowdata

\begin{document}
\begin{frame}[plain]
  \begin{tikzpicture}[remember picture,overlay]
    \begin{scope}[shift={(current page.north west)},
      every node/.style={anchor=north west,inner sep=0pt,outer sep=0pt,align=left}]
      \node[text width=146mm,font=\fontsize{17}{20}\selectfont\bfseries]
        at (7mm,-5mm) {Перенос VQDM на DiT: важны средние шаги};
      \node[text width=146mm,text=Subtle,font=\fontsize{9.5}{11.5}\selectfont]
        at (7mm,-15mm) {DiT-XL/2\quad ImageNet $256\times256$\quad 50 шагов DDPM\quad CFG 1,5};
      \node[text width=146mm,font=\fontsize{10.5}{12.5}\selectfont]
        at (7mm,-21mm) {\textbf{VQDM:} $32\to4{,}15$ бит/параметр\hfill
          \textbf{FID$\downarrow$:} $6{,}19\to15{,}73$\ \textcolor{Subtle}{(10\,000 изображений)}};
      \node[text width=96mm,font=\fontsize{10.5}{12.5}\selectfont\bfseries]
        at (7mm,-29mm) {Возврат исходной модели на 10 шагов};
      \node at (7mm,-35mm) {\begin{tikzpicture}[every node/.style={}]
        \begin{axis}[
          width=94mm,height=37mm,
          axis x line*=bottom,axis y line*=left,
          xmin=-.65,xmax=4.65,ymin=-3,ymax=49,
          xtick={0,1,2,3,4},
          xticklabels={0--9,10--19,20--29,30--39,40--49},
          ytick={0,10,20,30,40},
          xlabel={Шаги генерации с исходной моделью},
          ylabel={Снижение NMSE, \%},
          label style={font=\fontsize{9}{11}\selectfont},
          tick label style={font=\fontsize{9}{11}\selectfont},
          tick style={draw=Ink,line width=.4pt},
          axis line style={draw=Ink,line width=.45pt},
          ymajorgrids=true,grid style={draw=black!10,line width=.3pt},
          clip=false,
          /pgf/number format/use comma,
        ]
          \addplot[ybar,bar width=9mm,bar shift=0pt,fill=Muted,draw=none]
            table[x=position,y=gain]{\windowdata};
          \addplot[ybar,bar width=9mm,bar shift=0pt,fill=Accent,draw=none]
            coordinates {(2,36.488425)};
          \addplot[only marks,mark=none,
            error bars/.cd,y dir=both,y explicit,
            error bar style={draw=Ink,line width=.5pt},
            error mark options={draw=Ink,rotate=90,mark size=2pt,line width=.5pt}]
            table[x=position,y=gain,
              y error minus expr=\thisrow{gain}-\thisrow{low},
              y error plus expr=\thisrow{high}-\thisrow{gain}]{\windowdata};
          \addplot[only marks,mark=none,point meta=explicit,
            nodes near coords,
            every node near coord/.append style={
              font=\fontsize{10}{12}\selectfont\bfseries,
              text=Ink,anchor=south,inner sep=0pt,
              /pgf/number format/fixed,/pgf/number format/precision=1,
              /pgf/number format/fixed zerofill}]
            table[x=position,y expr=\thisrow{high}+2.4,meta=gain]{\windowdata};
        \end{axis}
      \end{tikzpicture}};
      \node[text width=43mm,font=\fontsize{9.5}{11.5}\selectfont]
        at (110mm,-29mm) {\setlength{\parskip}{0pt}\textbf{Проверка}\par\vspace{1mm}
          Исходная: 10 шагов.\par
          Квантованная: 40 шагов.\par\vspace{1.5mm}
          256 пар, общие шумы.\par
          Без сброса латента.};
      \node[text width=43mm,text=Accent,font=\fontsize{14}{16}\selectfont\bfseries]
        at (110mm,-58mm) {$-36{,}5\%$ NMSE};
      \node[text width=43mm,font=\fontsize{9.5}{11.5}\selectfont]
        at (110mm,-65mm) {при возврате шагов 20--29.};
      \node[text width=43mm,text=Subtle,font=\fontsize{9}{11}\selectfont]
        at (110mm,-74mm) {Диагностика,\par не улучшение FID.};
      \node[text width=96mm,text=Subtle,font=\fontsize{8.5}{10.5}\selectfont]
        at (7mm,-74mm) {NMSE: нормированная ошибка конечного латента\par
          относительно исходной модели. Интервалы: 95\%.};
      \node[text width=146mm,font=\fontsize{10.5}{12.5}\selectfont\bfseries]
        at (7mm,-84mm) {Идея: учитывать важность шагов при обучении кодбуков.};
    \end{scope}
  \end{tikzpicture}
\end{frame}
\end{document}
